% resumo em português
\begin{resumo}
    Segundo a NBR6028, o resumo deve ressaltar o objetivo, o método, os resultados e as conclusões do documento. A ordem e a extensão destes itens dependem do tipo de resumo (informativo ou indicativo) e do tratamento que cada item recebe no documento original. Descreva de maneira sucinta o tema do trabalho, explicando o que o motivou a realizá-lo e sobre o que será tratado. O importante é que alguém possa compreender do que trata o trabalho pela leitura deste resumo. No resumo não deve conter citações. As palavras-chave devem figurar logo abaixo do resumo, antecedidas da expressão ``Palavras-chave:'', separadas entre si por ponto e finalizadas também por ponto.
    
    \vspace{1em}
    \textbf{Palavras-chave:} modelo. LaTeX. PPGEE. IFPB.
\end{resumo}


% resumo em inglês
\begin{resumo}[Abstract]
 \begin{otherlanguage*}{english}
    Describe the theme of the work in a brief way, explaining what motivated you to accomplish it and about what it will be. The important is that someone can understand what the work is about, simply by reading this abstract.

   \vspace{1em}
   \textbf{Keywords:} template. LaTeX. PPGEE. IFPB.
 \end{otherlanguage*}
\end{resumo}